\subsection{多域融合的优势}
\label{sec:discussion-fusion}

第~\ref{sec:results-ablation}节中的消融结果表明，多域融合是 PMDF-Net 中最具影响力的单一架构选择。
移除时频分支和融合模块（无多域的 PMDF-Net）会使信噪比(SNR)相对于完整模型下降$2.1$dB，而仅使用时域操作会损失$2.6$dB，仅使用时频操作会损失$3.8$dB。这些差距并非微不足道；在激光超声(LU)信号处理中，$2.6$dB 的信噪比(SNR)提升直接转化为可检测缺陷尺寸的提升，否则这些缺陷将低于仪器分辨率极限。

根本原因在于表征的互补性。时域捕捉时序和瞬态波形形态，包括表征窄带缺陷的 abrupt phase shifts 和局部脉冲畸变。
频域编码频谱内容——能量在各谐波带、模态频率和共振特征上的分布，这些特征对瞬态事件精确时序具有不变性。
使得反射器到达时间位移微秒级的缺陷会在时域产生可测量的相位异常，而同一缺陷可能同时将频谱能量在各共振模态间转移，转移量在频域可被检测到。
单一视角无法保留足够信息来区分缺陷相关变化与噪声引起的伪影。

这种互补性对于缺陷检测尤为关键。许多具有物理意义的缺陷特征并不主导单一域；它们分散在两个域中。
降低界面刚度的分层会使共振峰轻微展宽，但不会显著改变主回波的到达时间。
相反，引入次级微回波的疲劳裂纹会以 leaving spectral features largely intact 的方式扰动时域波形。
PMDF-Net 中的融合模块学习对这些域特定特征进行加权和对齐，使任一域中的弱信号被另一域中的强信号放大，从而实现在任一单域配置下均会失败的检测决策。
因此，融合带来的$2.6$dB 增益不仅仅是一个信噪比(SNR)指标；它是可检测缺陷严重程度下限的实际改进。